% T17 source: computational_results.tex lines 1--59
The series below are organized by substantive question rather than by numerical
identifier: attainability of the bound and the role of phase ordering are
examined first, followed by parallelism and decomposition, and then by scale
robustness, extension to synthetic DAGs, and controlled placement of human
work. A numerical identifier denotes a frozen protocol, not the position of a
series in the logic of the exposition.

\subsection{Attainability of \texorpdfstring{$B_4$}{B4}}

The equality obtained for the illustrative scenario in EXP-00 is not a general
property of the model. In EXP-01, all 90 scenarios were solved to optimality,
and in 30 of them a strict gap $T^*>B_4$ was found; the maximum relative gap on
this grid was 18.333\%. The selected four-task counterexample on the EXP-01
grid is the diamond at $P=2$: $B_4=48.00$, whereas $T^*=49.25$ hours. The
additional 1.25 hours arise from a resource chain of repeated human-attention
phases that is absent from the three aggregate terms defining $B_4$.

EXP-03 directly tests the sufficiency of the aggregates. In 14 of the 24 pairs,
identical DAGs and identical values of $A$, $H$, $W_4$, $L_4$, $B_4$, $C$, and
$\gamma$ yielded different values of $T^*$. In the most pronounced selected
pair, with a common $B_4=36.0$, the two phase profiles yielded 40.5 and 37.75
hours. Thus, $B_4$ correctly provides a structural lower bound, whereas the
exact value is determined by the resource-augmented phase graph.

Figure~\ref{fig:resource-augmented-path} shows such a chain explicitly at one
EXP-05 point with $P=2$, $m=4$, and an agent share of overhead equal to $0.05$.
In this chain, the ordinary phase sequence and DAG dependencies are interleaved
with the service order of the single human resource and the task order on an
agent slot. Its length of $36.45$ hours therefore equals the certified $T^*$ and
exceeds the aggregate lower bound $B_4=31.80$ hours.

\begin{figure}[H]
  \centering
  \includegraphics[width=0.96\textwidth]{candidate_resource_augmented_path.png}
  \caption{Resource-augmented critical path: task sequence --- the phase order
  within a task, precedence --- a DAG edge, human order --- the order on the
  human resource, and slot order --- the task order on one agent slot. The
  labels $d1,d2$ refer to the two decomposed subtasks.}
  \label{fig:resource-augmented-path}
\end{figure}

\subsection{Configured Parallelism}

Across the 70 exact EXP-02 scenarios with $C=\gamma=1$ held constant,
increasing $P$ solely as available capacity did not worsen the optimum and
quickly produced a plateau determined by either the critical-path term or the
human-resource term. When instead an increase in $P$ was accompanied by an
increase in $C(P)$ or $\gamma(P)$, an interior minimum occurred at $P=2$ in all
four corresponding curves examined.
This is consistent with Propositions~\ref{prop:capacity-monotonicity}
and~\ref{prop:finite-optimum}: the deterioration is caused not by an additional
unused stream but by a change in the process configuration.

\begin{figure}[H]
  \centering
  \includegraphics[width=0.94\textwidth]{exp02_scaling_curve.png}
  \caption{Exact makespan as the number of agent streams varies. Constant
  overhead produces saturation, whereas increasing $C(P)$ or $\gamma(P)$
  produces an interior minimum on the grid examined.}
  \label{fig:exp02-scaling}
\end{figure}
